%==========================================================
\section{Introdução}\label{sec:introducao}

\subsection{O projeto \projeto}

O \projeto{} é um agente de viagens baseado em inteligência artificial, 
pensado, numa primeira fase, para jovens com orçamentos limitados
mas posteriormente para qualquer utilizador que queira simplificar o processo de planeamento de viagens. 
Ainda pretendemos avaliar o potencial do \projeto{} em atender às novas necessidades de agências de viagens já 
inseridas no mercado. As plataformas atuais 
estão fragmentadas: obrigam o utilizador a escolher primeiro um destino e a 
pesquisar à parte voos, alojamento e atividades. Pretendemos combinar 
estas opções automaticamente, a partir do orçamento, das datas, dos interesses 
e da flexibilidade de cada pessoa, e propor um itinerário personalizado e 
acessível.

Antes de decidir o que construir, é preciso perceber como esse trabalho já é
feito hoje. Os jovens já planeiam viagens baratas, com abas abertas, folhas
de cálculo, capturas de ecrã e grupos de mensagens. Os agentes de viagens já
fazem, ao balcão, parte do que o \projeto{} quer automatizar, e são ao mesmo tempo
potenciais clientes: interessa perceber que necessidades novas têm e onde uma ferramenta
como o \projeto{} os poderia apoiar. Esta entrega estuda o trabalho atual destes dois
grupos: o jovem que planeia a viagem e o agente que a planeia por profissão.

\subsection{Objetivos da entrega}

A Tabela~\ref{tab:mapa-entrega} resume os elementos pedidos no enunciado da Entrega~1 e onde estão neste relatório.

\begin{table}[!htbp]
\centering\small
\begin{tabularx}{\textwidth}{@{}p{0.6cm}Lp{5.2cm}@{}}
\toprule
\textbf{N.º} & \textbf{Elemento pedido no enunciado} & \textbf{Onde está}\\
\midrule
1 & Notas de campo e conjunto consolidado de modelos de trabalho (fluxo, sequência e pelo menos um de artefacto, cultural ou físico) & Secções~\ref{sec:recolha}, \ref{sec:modelos} e~\ref{sec:consolidacao}; Apêndice~\ref{anx:notas}\\
2 & Diagrama de afinidade que consolida as sessões & Secção~\ref{sec:afinidade}  \\
3 & Protótipo de baixa fidelidade ou \textit{storyboard} & Secção~\ref{sec:prototipo} \\
4 & Avaliação heurística e percurso cognitivo, por quem não construiu o protótipo & Secção~\ref{sec:avaliacao} \\
5 & Reflexão de uma página: onde e porquê as duas avaliações discordaram & Secção~\ref{sec:reflexao}\\
6 & Apêndice de uso de IA, com um exemplo do que a IA errou & Apêndice~\ref{anx:ia}\\
\bottomrule
\end{tabularx}
\caption{Correspondência entre o enunciado da Entrega~1 e este relatório.}
\label{tab:mapa-entrega}
\end{table}

\subsection{Estrutura do documento}

A Secção~\ref{sec:planeamento} descreve o planeamento da investigação contextual. 
A Secção~\ref{sec:recolha} apresenta as sessões realizadas e a forma como foram registadas. 
A Secção~\ref{sec:modelos} apresenta os modelos de trabalho individuais e a Secção~\ref{sec:consolidacao} 
a sua consolidação. As Secções~\ref{sec:afinidade} a~\ref{sec:avaliacao} tratam do diagrama de afinidade, 
do protótipo e da sua avaliação. A Secção~\ref{sec:reflexao} contém a reflexão crítica e a 
Secção~\ref{sec:conclusao} fecha com as conclusões e os próximos passos. Os anexos reúnem as notas de campo 
completas, os modelos em tamanho integral, o apêndice de uso de IA, a organização da equipa e as hipóteses de 
produto que os dados sustentam.
